\section*{Capítulo 2}

\solucion{2.1} Simple: $5{,}000(1+0.06\times4)=6{,}200.00$. Compuesto: $5{,}000\times1.06^4=5{,}000\times1.262477=6{,}312.38$. La diferencia de US\$ 112.38 son los intereses ganados sobre los intereses reinvertidos.

\solucion{2.2} $\VP=500/1.0472^2=500/1.096628=455.94$ millones de dólares.

\solucion{2.3} Factor a 91 días: $1+0.039\times91/360=1.009858$. Efectiva anual: $1.009858^{360/91}-1=3.957\,\%$. Es la tasa que el informe redondea al construir el NDF.

\solucion{2.4} $C=1{,}000\times\dfrac{0.01}{1-1.01^{-12}}=1{,}000\times\dfrac{0.01}{1-0.887449}=88.85$. Primera cuota: interés $=1{,}000\times0.01=10.00$; amortización $=88.85-10.00=78.85$; saldo $=921.15$.

\solucion{2.5} $DF(1)=1/1.04=0.961538$. $DF(2)=\dfrac{1-0.045\times0.961538}{1.045}=\dfrac{0.956731}{1.045}=0.915532$. Tasa cupón cero: $z(2)=0.915532^{-1/2}-1=4.511\,\%$, ligeramente mayor que la tasa swap (4.50\,\%) porque la curva es creciente.

\solucion{2.6} El factor de descuento traduce un flujo futuro a su valor hoy \emph{en la misma moneda}, usando el costo de oportunidad de esa moneda. Las tasas en soles y en dólares difieren por inflación esperada, política monetaria y riesgo país. Descontar soles con tasas en dólares produciría un valor en soles que no corresponde a ningún precio de mercado; el error se trasladaría al nocional y a la tasa del CCS.

\solucion{2.7} No. 1\,\% mensual compuesto da una TEA de $1.01^{12}-1=12.68\,\%$. La tasa mensual equivalente a una TEA de 12\,\% es $1.12^{1/12}-1=0.9489\,\%$.
